\section{Materiais e Métodos}
\label{sec:material_metodos}

A metodologia compreendeu a construção do nariz eletrônico, o delineamento
experimental, a aquisição dos sinais do sistema solo--planta, o
pré-processamento das séries temporais e a classificação das coletas a vácuo.

\subsection{Sistema de nariz eletrônico}
\label{subsec:sistema_nariz}

O nariz eletrônico foi desenvolvido para registrar a assinatura de compostos
voláteis.
Foi composto por microcontrolador ESP32, seis sensores de gases MQ, sensor de
umidade do solo, sensor BMP280 de temperatura e pressão, bomba de vácuo,
válvulas, tubo de PVC acoplado ao vaso, mangueiras, câmara sensorial de vidro
e estrutura externa de MDF. O ESP32 registrou os sinais e controlou a bomba
e as válvulas; uma interface web acessada por Wi-Fi permitiu iniciar e
acompanhar as aquisições.
O BMP280 forneceu leituras de temperatura e pressão, mas não mediu a umidade
relativa do ar na câmara sensorial. A leitura do sensor de solo caracterizou
um índice operacional de umidade do substrato; ela não foi usada como medida
de umidade relativa do gás apresentado aos sensores MQ.

\begin{figure}[H]
    \centering
    \includegraphics[width=0.30\textwidth]{figures/completo.jpeg}
    \caption{Estrutura externa do nariz eletrônico.}
    \label{fig:estrutura_nariz_eletronico}
\end{figure}

\subsubsection{Matriz de sensores de gases}
\label{subsubsec:matriz_sensores}

A matriz foi constituída por MQ-2, MQ-3, MQ-7, MQ-8, MQ-135 e MQ-138,
posicionados na câmara de vidro sob o mesmo fluxo de amostra. Esses sensores
semicondutores de óxido metálico apresentam sensibilidades parcialmente
sobrepostas. Por isso, a classificação utilizou o padrão conjunto das seis
respostas, sem atribuir a um sensor isolado a identificação química de um
composto. As curvas de sensibilidade dos fabricantes foram empregadas apenas
para descrever as respostas nominais dos modelos MQ.

\begin{figure}[H]
    \centering
    \includegraphics[width=0.30\textwidth]{figures/matriz.png}
    \caption{Matriz dos seis sensores de gases na câmara sensorial.}
    \label{fig:matriz_sensores}
\end{figure}

\subsection{Delineamento experimental e material vegetal}
\label{subsec:delineamento}

O experimento foi conduzido na Estação de Pesquisas Agronema, em Grandes
Rios, Paraná, Brasil, nas coordenadas $24{,}186^{\circ}$ S e
$51{,}399^{\circ}$ O. Foram utilizados dez vasos com plantas de soja,
distribuídos em dois grupos experimentais de cinco vasos cada: controle e
inoculado. Os vasos constituíram as unidades biológicas do experimento.
As leituras sucessivas de um mesmo vaso e as coletas repetidas ao longo dos
dias não foram tratadas como réplicas biológicas independentes.

As aquisições começaram aproximadamente 35 dias após a inoculação. A base
utilizada para a análise com coleta a vácuo contém 36 coletas.




\subsection{Vasos controle e vasos inoculados com \textit{Heterodera glycines}}
\label{subsec:classes_experimentais}

Conforme o delineamento experimental, os vasos 1 a 5 formaram o grupo
controle e os vasos 6 a 10 formaram o grupo inoculado. Os vasos inoculados
receberam aproximadamente 100 mil juvenis de segundo estádio (J2) de
\textit{H. glycines} por vaso.
Foi empregada uma única dose de inoculação, de modo que o delineamento não
permite avaliar resposta a diferentes densidades do nematoide. Os rótulos
usados na classificação representam a alocação experimental (inoculado ou
controle); contagens finais de cistos, ovos ou juvenis por vaso não foram
incorporadas como padrão de referência nesta análise. Assim, a classe
inoculada não equivale à confirmação individual de infestação.

A Figura \ref{fig:vasos_experimento} apresenta os vasos
utilizados no experimento, divididos entre os grupos
controle e inoculado.

\begin{figure}[H]
    \centering
    \includegraphics[width=0.55\textwidth]{figures/vasos_experimento.png}
    \caption{Vasos de soja utilizados no experimento,
    divididos entre os grupos controle e inoculado com
    \textit{H. glycines}.}
    \label{fig:vasos_experimento}

    \vspace{0.1cm}
    {\footnotesize Fonte: Elaborado pelos autores (2026).}
\end{figure}

Na base de classificação, a \textbf{classe 0} identifica
o grupo inoculado e a \textbf{classe 1} identifica o grupo controle.
Assim, ao longo do artigo, o recall da classe inoculada corresponde
ao recall da classe 0. A designação de classe não constitui, por si,
confirmação nematológica da condição de cada vaso.


\subsection{Protocolo de aquisição de dados por coleta a vácuo}
\label{subsec:protocolo_aquisicao}

O tubo de PVC foi acoplado ao vaso para delimitar o espaço de amostragem
(Figura~\ref{fig:vasos_experimentos}). A bomba conduziu o ar do
\textit{headspace} pela tubulação até a câmara dos sensores. Uma execução
completa do ciclo, registrada sob um identificador próprio de vaso e dia, foi
definida como uma coleta. Em cada leitura, o ESP32 registrou o tempo decorrido,
os seis sinais MQ, a leitura do sensor de solo, a temperatura e a pressão.
As coletas a vácuo analisadas ocorreram nos dias experimentais 11, 12, 13,
18 e 19. As coletas posteriores realizadas sob outro protocolo não
integraram esta análise. Não se atribuiu à leitura de pressão uma relação
exclusiva com a composição dos gases, pois ela também pode variar com a
operação da bomba, a tubulação e a vedação do sistema.

\begin{figure}[H]
    \centering
    \includegraphics[width=0.55\textwidth]{figures/vasos_experimentos.png}
    \caption{Protótipo posicionado sobre o vaso de soja durante a coleta a vácuo.}
    \label{fig:vasos_experimentos}
    \vspace{0.1cm}
    {\footnotesize Fonte: Elaborado pelos autores (2026).}
\end{figure}

A Figura~\ref{fig:ciclo_coleta} apresenta as etapas de operação do sistema
com e sem controle de pressão. Neste estudo, foram analisadas as coletas
realizadas com controle de pressão.

\begin{figure}[H]
    \centering
    \includegraphics[width=\textwidth]{figures/fluxo_en.png}
    \caption{Fluxograma dos ciclos de coleta com e sem controle de pressão
    e esquema do sistema de amostragem.}
    \label{fig:ciclo_coleta}
    \vspace{0.1cm}
    {\footnotesize Fonte: Elaborado pelos autores (2026).}
\end{figure}

O ciclo de referência, ilustrado na Figura~\ref{fig:ciclo_coleta},
compreendeu aquecimento dos sensores, aproximadamente entre 0 e 180~s; uma
transição operacional até cerca de 200~s; amostragem do \textit{headspace},
aproximadamente entre 200 e 800~s; e desaturação ou recuperação,
aproximadamente entre 800 e 900~s. Esses intervalos descrevem o protocolo de
referência, não limites rígidos aplicados a todas as séries: a duração efetiva
das coletas variou, e parte dos registros terminou antes de 900~s. O controle
da bomba e das válvulas foi realizado pelo ESP32 durante o ciclo.

A condição denominada solo seco ou úmido foi definida pelo índice operacional
\texttt{Soil\_indice\_0\_1}, obtido pela normalização mínimo--máximo da leitura
do sensor de umidade do solo na base consolidada, conforme a
Equação~\eqref{eq:indice_solo}:
\begin{equation}
    I_{\mathrm{solo}} =
    \frac{\mathrm{Soil}-\min(\mathrm{Soil})}
         {\max(\mathrm{Soil})-\min(\mathrm{Soil})},
    \label{eq:indice_solo}
\end{equation}
em que $\mathrm{Soil}$ é a leitura bruta do sensor de umidade do solo, e
$\min(\mathrm{Soil})$ e $\max(\mathrm{Soil})$ são, respectivamente, os valores
mínimo e máximo dessa leitura na base consolidada.
O mínimo e o máximo foram calculados sobre a base consolidada antes da
separação entre treino e teste; por isso, as leituras de teste participaram
da definição dessa escala. Não se trata de uma normalização ajustada
exclusivamente no treinamento.

Leituras com $I_{\mathrm{solo}}\leq 0{,}4$ foram denominadas \emph{solo seco};
leituras com $I_{\mathrm{solo}}>0{,}4$, \emph{solo úmido}. No conjunto de
teste, todas as leituras de cada coleta ficaram do mesmo lado desse limiar,
permitindo classificar também a coleta inteira como seca ou úmida. O índice não
representa porcentagem calibrada de água no solo. Os termos \emph{seco} e
\emph{úmido} designam apenas duas faixas operacionais do índice, separadas
por um limiar de 0,4. Não foi estabelecida uma relação causal entre essa
classificação e o desempenho dos modelos.

\subsection{Pré-processamento dos sinais e janelamento}
\label{subsec:preprocessamento}

Cada planilha da base original representou uma coleta. As leituras foram
associadas ao identificador da coleta, ao dia, ao vaso e à classe. Na
consolidação, foram removidos 15\% dos registros iniciais e 15\% dos
registros finais de cada série para reduzir a contribuição dos transientes
de início e fim do ciclo. Esse recorte proporcional não equivale a impor
os mesmos limites de tempo a todas as coletas: permaneceu a faixa central
de 70\% de cada série, que pode conter trechos fora da amostragem nominal
de 200 a 800~s. As colunas necessárias foram
convertidas para formato numérico, e registros sem valores válidos nessas
colunas foram descartados.
Portanto, a faixa retida não corresponde necessariamente apenas à fase
nominal de exposição ao \textit{headspace}; esse recorte foi aplicado de
forma igual a todas as coletas.

Para remover transientes de pressão, calculou-se
$\Delta P_t=|P_t-P_{t-1}|$ dentro de cada coleta. Quando
$\Delta P_t\geq 0{,}10$~kPa, foram excluídas a leitura do evento e até
30 leituras anteriores e 30 posteriores. Na versão estrita do filtro,
também foram removidas leituras fora da faixa definida pela mediana global
da pressão $\pm 0{,}50$~kPa, calculada sobre a base consolidada. Após a
filtragem estrita, a base de coletas a vácuo continha 63.492 leituras de
36 coletas. Não foi aplicada média móvel nem outra suavização temporal
aos sinais usados pelos classificadores.
A mediana e os limites da faixa estrita foram calculados antes da divisão
entre treino e teste, de modo que o conjunto de teste participou da
definição desse filtro operacional. Os limiares de $0{,}10$ e $0{,}50$~kPa
foram empregados neste processamento; seu efeito sobre as métricas não
foi isolado por uma análise de sensibilidade nesta seção.

A qualidade dos sinais também foi examinada por coleta. Um registro com
ruído não implica descartar as demais coletas do mesmo vaso: na auditoria,
a coleta inoculada rotulada como ``vaso 1'' no dia 19 apresentou oscilações
multissensor acentuadas, enquanto as coletas com o mesmo identificador nos
dias 12 e 18 não apresentaram esse padrão. Uma coleta controle do dia 19
também foi sinalizada. Essas observações motivam uma análise de sensibilidade
com exclusões definidas por qualidade do sinal, aplicada com o mesmo critério
às duas classes. Os totais e resultados principais relatados aqui ainda
correspondem às 36 coletas, sem essas exclusões.

Para a modelagem principal, as entradas foram seis sinais MQ corrigidos
estatisticamente, o índice de solo, a temperatura em $^{\circ}$C e a
pressão em kPa. Para cada MQ, ajustou-se somente com as coletas de
treinamento um \texttt{HuberRegressor} que estimou a componente associada
ao índice de solo, à temperatura e à pressão. Essa componente foi
centralizada pela média das previsões no treinamento e subtraída do
sinal original, conforme a Equação~\eqref{eq:compensacao_mq}:
\begin{equation}
    \label{eq:compensacao_mq}
    MQ_{s,t}^{\mathrm{corr}} = MQ_{s,t} -
    \left(\widehat{MQ}_{s,t}^{\mathrm{amb}}-
    \overline{\widehat{MQ}_{s}^{\mathrm{amb}}}_{\mathrm{treino}}\right).
\end{equation}
Nessa equação, $s$ identifica o sensor MQ, $t$ identifica a leitura, e a
barra representa a média das previsões nas leituras de treinamento. Os
``MQ corrigidos'' são, portanto, sinais com compensação \emph{estatística}
da associação linear estimada com essas três variáveis. Eles não resultam
de calibração física para concentração de gases ou umidade relativa do ar.
As seis entradas MQ, o índice de solo, a temperatura e a pressão totalizaram
nove canais por leitura.

Para a TCN, as leituras foram ordenadas por coleta e tempo e segmentadas
em janelas de 128 leituras, com passo de 64 leituras e sobreposição nominal
de 50\%, sem atravessar limites entre coletas. Quando necessário, uma
janela adicional foi alinhada ao fim da coleta para incluir as últimas
leituras; sua sobreposição podia superar 50\%. Cada janela teve formato de
$9\times128$ (canais $\times$ leituras). As nove entradas foram padronizadas
com parâmetros ajustados somente nas coletas de treinamento. Para coletas
com menos de 128 leituras, a janela foi completada com repetição do último
valor disponível. O classificador arbóreo recebeu as entradas por leitura,
sem essa padronização adicional.


\subsection{Modelos de aprendizado de máquina}
\label{subsec:modelos}

Foram comparados um classificador de árvores extremamente aleatorizadas,
\emph{Extra Trees} (ET) \cite{geurts2006extremely}, e uma rede convolucional
temporal residual dilatada (TCN) \cite{bai2018tcn}. A implementação local
do modelo arbóreo usa \texttt{ExtraTreesClassifier}, com 900 árvores,
\texttt{max\_features="sqrt"}, mínimo de dez amostras por folha,
\texttt{bootstrap=False}, sem pesos de classe e semente 42.

A TCN implementada neste trabalho usou três blocos convolucionais
temporais residuais, com 24, 24 e 32 canais, convoluções de núcleo 3 e
dilatações 1, 2 e 4. Cada bloco contém duas convoluções unidimensionais
seguidas por ativação ReLU; o caminho residual usa uma convolução de núcleo
1 quando o número de canais muda. As convoluções empregam preenchimento
de $2d$ posições para dilatação $d$, com recorte da saída ao comprimento
da entrada após cada bloco. Não há camada de normalização dentro dos blocos.
A saída dos blocos passou por agregação média global
no tempo, \emph{dropout} de 0,20 e camada linear de duas classes. O
treinamento utilizou 60 épocas, lotes de 32 janelas, otimizador AdamW com
taxa de aprendizado $10^{-3}$ e \emph{weight decay} de $10^{-4}$, e
entropia cruzada com pesos inversamente proporcionais à frequência das
classes nas janelas de treinamento. A semente aleatória foi 42. A rede
produziu probabilidades por janela; para avaliar cada leitura, foi usada
a média das probabilidades das janelas que a cobriam, com limiar 0,5 para
a classe 1. A arquitetura não inclui mecanismo de atenção nem método de
explicabilidade integrado. O número de épocas foi fixado em 60 para esta
comparação, sem seleção por validação interna. Cada modelo foi treinado uma
vez com a semente indicada.

Cada classificador foi ajustado uma vez com todas as coletas de treinamento
sob pressão forçada. Solo seco e solo úmido foram dois recortes do mesmo
conjunto de teste, não dois treinamentos independentes.

\subsection{Estratégia de treinamento e teste}
\label{subsec:treinamento_teste}

A divisão foi feita por \emph{coleta}, dentro de cada classe, com
proporção nominal de 70\% para treinamento e 30\% para teste e semente 42.
Todas as leituras e janelas de uma coleta permaneceram em um único conjunto.
Das 36 coletas, 24 foram destinadas ao treinamento (41.377 leituras: 11
coletas da classe 0 e 13 da classe 1) e 12 ao teste (22.115 leituras:
seis coletas de cada classe).
Os hiperparâmetros descritos acima foram usados nessa única partição;
não foi realizada validação cruzada agrupada para escolhê-los no conjunto
de treinamento.

Essa separação evita que janelas da mesma aquisição apareçam nos dois
conjuntos. Entretanto, alguns vasos aparecem em ambos porque foram medidos
em dias distintos. Assim, o teste avalia novas \emph{coletas}, não vasos
biologicamente independentes, e as métricas não demonstram generalização
para plantas ou vasos nunca vistos pelo modelo.
As 63.492 leituras derivam de 36 coletas feitas em dez vasos; elas não
correspondem a 63.492 réplicas biológicas. A correlação entre leituras de
uma mesma coleta e entre coletas de um mesmo vaso limita a interpretação
da precisão aparente das métricas por leitura.

No teste, o recorte seco incluiu 16.692 leituras de oito coletas: 6.507
leituras da classe 0 e 10.185 da classe 1. O recorte úmido incluiu 5.423
leituras de quatro coletas: 4.152 da classe 0 e 1.271 da classe 1. Essa
distribuição desigual motivou o relato da acurácia balanceada junto da
acurácia convencional.

\subsection{Avaliação exploratória agrupada por vaso}
\label{subsec:validacao_vaso}

Como análise de sensibilidade à escolha da unidade de separação, cinco
configurações de entradas previamente selecionadas foram avaliadas com
Extra Trees, sem novo ajuste de hiperparâmetros: (i) MQ2, MQ135 e MQ138
com variáveis ambientais; (ii) MQ2, MQ3, MQ8, MQ135 e MQ138; (iii) os seis
MQ; (iv) os seis MQ com variáveis ambientais; e (v) MQ7 com variáveis
ambientais. Nessa análise, os sinais MQ foram usados em sua forma bruta.
Cada coleta recebeu uma previsão única, obtida pela média das
probabilidades de suas leituras, com limiar de 0,5 para a classe 1. Para
identificar unidades biológicas na base, combinou-se a classe experimental
e o número registrado do vaso; essa chave técnica evitou a colisão do
identificador ``vaso 1'', presente nas duas classes do arquivo analisado.

Além da partição histórica por coleta, foram examinados dois esquemas sem
compartilhar a chave de vaso entre treino e teste: uma divisão estratificada
por classe com sete vasos no treinamento e três no teste (semente 42) e dez
rodadas \textit{leave-one-vase-out}, treinando em nove vasos e avaliando no
vaso restante. Na segunda estratégia, a acurácia balanceada foi calculada
depois de reunir as previsões das 36 coletas, pois cada rodada de teste
continha somente uma classe. Nesse procedimento, nenhum vaso foi eliminado
da avaliação: cada um integrou o teste em uma das rodadas. Essa análise foi
restrita ao Extra Trees e
às configurações selecionadas anteriormente com os mesmos dados; seus
resultados são exploratórios e não constituem validação externa nem
estimativa independente da generalização da TCN.

\subsection{Métricas de desempenho}
\label{subsec:metricas}

Foram calculadas acurácia, acurácia balanceada, F1 macro e matrizes de
confusão no conjunto de teste. A acurácia é a fração de leituras
classificadas corretamente; a acurácia balanceada é a média dos recalls
das classes 0 e 1; e a F1 macro é a média simples das medidas F1 das duas
classes. O recall da classe inoculada corresponde à sensibilidade
para a classe 0, ou seja, à fração de leituras de vasos inoculados
corretamente identificadas. As métricas foram calculadas separadamente
para os recortes seco e úmido do mesmo teste. Os denominadores dessas
métricas são leituras, não coletas ou vasos; coletas mais longas podem,
portanto, receber maior peso. Não foram estimados intervalos de confiança
por reamostragem de vasos nesta avaliação principal.

\subsection{Análise complementar das entradas}
\label{subsec:analise_entradas}

Para examinar a influência das variáveis ambientais, os dois classificadores
foram ajustados novamente com a mesma partição por coleta em dois cenários:
seis sinais MQ brutos e os mesmos seis sinais acrescidos do índice de solo,
da temperatura e da pressão. Os MQ brutos foram usados em ambos os cenários
para que a comparação não misturasse a inclusão de variáveis ambientais
com a compensação estatística dos MQ utilizada na análise principal. Por
essa razão, as métricas desta análise complementar não devem ser comparadas
diretamente com as dos modelos principais de MQ corrigidos.

A importância por permutação foi calculada sobre as 22.115 leituras do
teste completo como a diferença, em pontos percentuais, entre a acurácia
balanceada original e aquela obtida após embaralhar os valores de uma
entrada entre leituras, mantendo o modelo treinado e as demais entradas
fixas \cite{fisher2019modelreliance}. Foram feitas três permutações por
entrada em cada cenário. Na TCN, as janelas foram reconstruídas após a
perturbação da série. Esse procedimento mede dependência preditiva sob
permutação por leitura; ele rompe a ordem temporal da entrada perturbada
e não identifica uma contribuição física ou causal do sensor. A pequena
quantidade de repetições também limita a avaliação da estabilidade da
ordenação das variáveis.
